\section{复杂性风险：我害怕 Manus 变得复杂}

前面讲制造业式交付和基础设施，这一章回到产品哲学。访谈最后，Peak 说自己最害怕 Manus 失去特色，内部最害怕 Manus 变得复杂。这个判断非常重要。很多产品在增长中会自然增加功能，以满足更多用户、更多场景、更多客户；但每增加一个东西，都会稀释其他东西。

\begin{figure}[H]
\centering
\includegraphics[width=0.96\textwidth]{complexity-risk.png}
\caption{复杂性风险：增长、功能、稀释、克制和飞轮。}
\end{figure}

\begin{knowledgebox}{读图：复杂性来自增长的吸引力}
图中从增长到功能，再到稀释和克制。产品增长会诱导团队加更多能力，但功能越多，用户越难理解核心价值。Manus 的挑战是既保持克制，又不牺牲持续增长。
\end{knowledgebox}

\subsection{不是怕竞争，而是怕失去独特价值}

Peak 认为，和大厂竞争本身不是最可怕的；当你担心被断供或被抄时，产品已经到了很好的状态。更大的风险是失去独有价值，用户不再知道为什么选择你。这也是一种竞争，但不是外部竞争先杀死你，而是内部复杂性先稀释你。

\begin{importantbox}{产品独特性的底线}
一个产品可以扩展，但不能让用户忘记它最核心解决什么问题。增长不能以失去心智为代价。
\end{importantbox}

\subsection{本章小结}

Manus 的复杂性焦虑，是所有 Agent 公司都会面对的问题：越成功越容易加功能，越加功能越可能失去清晰心智。

